%==============================================================================
%  3. Proposed LLM-Enhanced DMDE
%==============================================================================

\section{Proposed LLM-Enhanced DMDE}
\label{sec:method}

\subsection{Overall Framework}
\label{subsec:framework}

% 建议配一张总体框架图：figures/fig_framework.pdf
\begin{figure}[htbp]
  \centering
  % \includegraphics[width=0.9\linewidth]{fig_framework.pdf}
  \fbox{\parbox[c][4cm][c]{0.9\linewidth}{\centering TODO: framework figure \\
    \small (the LLM acts as an external controller, decoupled from the DE loop)}}
  \caption{Overall architecture of the proposed \llmdmde{} framework. The LLM
           interacts with the evolutionary loop through a single search-control
           module, while the mutation, crossover and selection operators remain
           unchanged.}
  \label{fig:framework}
\end{figure}

As shown in \figref{fig:framework}, the framework keeps the DE main loop intact
and attaches a single LLM-driven module to it: \searchctrl{}, invoked every $I$
generations to re-tune the three decoupled search mechanisms. In vanilla DMDE
the scale factor \Ffactor{} and the extinction rate \GMR{} are derived from
\CR{} by formulas~3-11 and~3-12, so the three can only move together. Here the
chain is broken: \CR{}, \Ffactor{} and \GMR{} become three independent
channels, and the LLM may hold or change each one on its own.

% 待补：说明 LLM 调用预算（interval=50）、失败回退策略
% （调用失败或输出不可解析时退回公式默认行为），以及为什么这样设计是可复现的
\todo{report the query budget ($I = 50$), the fallback strategy when a call
      fails or returns unparsable output (degrade to the DMDE formulas), and
      why the design is reproducible.}

\subsection{Action-Space Control Conditions}
\label{subsec:action-space}

To isolate the effect of the action space from the effect of the observations,
three conditions are defined that share the same observation set, the same
trigger mechanism and the same freeze guards, and differ only in what the model
is allowed to control:

\begin{itemize}
  \item \textbf{Decoupled}: the model controls \CR{}, \Ffactor{} and \GMR{}
        independently (the main condition).
  \item \textbf{Coupled}: the model controls \CR{} only; \Ffactor{} and
        \GMR{} follow from formulas~3-11 and~3-12 exactly as in vanilla DMDE.
        This reproduces the original coupling while keeping the same controller.
  \item \textbf{CR-locked}: \CR{} is pinned to an offline-selected constant
        and the model controls \Ffactor{} and \GMR{} only. This measures the
        residual effect of the two channels once the dominant one is removed.
\end{itemize}

A deliberate constraint of this design is that the model never sees a channel it
cannot move: under the coupled condition the candidate list for \Ffactor{} is
omitted from the prompt, and no freeze notice is injected for \GMR{} or
\Ffactor{}, because telling the model to hold a channel it does not own would
be self-contradictory.

\subsection{LLM-Guided Search Control}
\label{subsec:searchctrl}

\searchctrl{} re-tunes \CR{}, \Ffactor{} and \GMR{} every $I$ generations.
Unlike a one-shot prompt, the controller is fed back the outcome of its previous
decision: the parameter values in force, the number of generations elapsed, the
relative fitness change, the change in population diversity, and the constraint
feasibility readings. The decision loop is therefore closed, and the prompt
exposes the parameters actually used in each recent interval.

\begin{algorithm}[htbp]
  \caption{\searchctrl{}}
  \label{alg:searchctrl}
  \begin{algorithmic}[1]
    \State \textbf{Input:} current generation $g$, interval $I$, state
           features and the stage history of the last intervals
    \State \textbf{Output:} \CR{}, \Ffactor{} and \GMR{} for the next
           interval (each defaults to \textsc{hold})
    \State \textbf{if} $g \bmod I \neq 0$ \textbf{then} keep current values
    \State build prompt: state features, stage history, trigger reason,
           shadow contrast, freeze flags
    \State $(\CR^{\star}, \Ffactor^{\star}, \GMR^{\star}) \gets$ LLM(prompt)
    \State clip each value that changed to its admissible candidate set
    \State \Return $(\CR, \Ffactor, \GMR)$
  \end{algorithmic}
\end{algorithm}

The default action for every channel is \textsc{hold}: the model is asked first
whether a parameter should change at all, and only if it answers yes is it asked
for a value. Two further safeguards keep the loop from fitting noise. First,
being consulted is not treated as evidence in itself -- the controller also
consults the model on a fallback timer. Second, each channel carries a freeze
guard: a channel that has produced no measurable effect over several intervals
is frozen and its output forced to \textsc{hold}. A shadow population evolved
from the same starting point with a fixed \CR{} is additionally maintained; the
difference between the two populations over the same interval is the quantity
the model is asked to beat, and a difference below the attribution threshold is
reported as indistinguishable from the baseline.

% 待补：说明取值范围的裁剪策略与默认回退值，并给出 prompt 模板的要点
\todo{state the clipping range, the default fallback value, and the key points
      of the prompt template.}
